\section{Introduction}
    Durante o processo de desenvolvimento de jogos em que
    existe algum tipo de combate com Non-Playable Caracters 
    (NPC), os desenvolvedores buscam criar combates que 
    entretenham, divirtam e mantenham o jogadores engajados, 
    sem tornar os repetitivos. Como uma alternativa aos
    os métodos heuristicos tradicionais, a ultilização de 
    machine-learning se mostra uma altenativa promissora 
    para Dynamic Difficulty Adjustment (DDA) dos combates 
    em jogos.
    
    Para tanto, foram analizados diversos artigos 
    \cite{mortazavi:24} nos quais são abordados diversas 
    técnicas de implementação da mecânicas de combate utilizado
    machine-learning como, 
    Belgian Artificial Intelligence (BAI) \cite{mi:22}, 
    AI agent using deep reinforcement learning with self-play \cite{kim:20},
    Dynamic Difficulty Adjustment (DDA) \cite{moon:22},
    Reinforcement Learning (RL) \cite{pagalyte:20}
    entre outras em diversos gêneros de jogos.
    
    

    